\documentclass[12pt,reqno]{amsart}
\usepackage{/Users/johnmyers/GoogleDrive/tex/handout,/Users/johnmyers/GoogleDrive/tex/customcommands,polynom}
\usepackage{caption}

\hdr{MAT 240}{\S16.5 {\small Integrals in cylindrical and spherical coordinates, part 2}}

\begin{document}

\textbf{Suggested problems:} 1-41 odd

\bigskip
\hrule

\bigskip
\prob Using the diagram below, write down equations that relate the spherical coordinates $(\rho,\theta,\phi)$ of a point to its cartesian (rectangular) coordinates $(x,y,z)$.

\medskip
\includegraphics[scale=1.25]{1.pdf}


\bigskip
\prob 

\begin{enumerate}
\item Plot the point with spherical coordinates $(2,\pi/4,\pi/3)$ and find its cartesian coordinates.\vfill
\item Find the spherical coordinates of the point with cartesian coordinates $(0,2\sqrt{3},-2)$.\vfill
\end{enumerate}

\prob

\begin{enumerate}
\item Describe the equation whose equation in spherical coordinates is $\rho=1$.\vfill
\item Give an equation in spherical coordinates for the cone $z = \sqrt{x^2+y^2}$.\vfill
\item Give an equation in spherical coordinates for the plane $z=10$.
\end{enumerate}

\vfill
\newpage
\prob Using the diagram below, express the volume element $dV$ in spherical coordinates.

\medskip
\includegraphics[scale=1.25]{pic2.pdf}

\bigskip
\prob For each object below, set up a triple integral of a function $f$ over the object.

\medskip
\begin{enumerate}
\item \textcolor{white}{blank} \\
\includegraphics[scale=1.25]{pic6.pdf} \vfill
\item A piece of a sphere, angle at the center is $\pi/3$. \\

\includegraphics[scale=1.25]{2.pdf} \\ \vfill
\item \textcolor{white}{blank} \\
\includegraphics[scale=1.25]{pic7.pdf}
\end{enumerate}

\vfill
\newpage
\prob Evaluate
	\[\int_W \sin{(\phi)} \ dV
	\]
where $W$ is the region $0 \leq \theta \leq 2 \pi$, $0 \leq \phi \leq \pi/4$, and $1 \leq \rho \leq 2$.

\vfill
\prob Compute the volume of the solid pictured below.

\medskip
\includegraphics[scale=1.25]{3.pdf}

\vfill
\prob Evaluate
	\[\int_W (x^2+y^2) \ dV,
	\]
where $W$ lies between the spheres $x^2+y^2+z^2=4$ and $x^2+y^2+z^2=9$.

\vfill
\end{document}